\chapter*{Conclusiones y Recomendaciones}
\addcontentsline{toc}{chapter}{Conclusiones y Recomendaciones}

A continuación se presentan las conclusiones y recomendaciones asociadas al proyecto materia del presente trabajo.

% Las conclusiones tampoco cuentan como capítulo

\section*{Conclusiones}

\begin{itemize}
    \item En este documento se ha presentado un producto de datos encaminado a calcular y determinar diferentes elementos de valor para la industria  de investigación de mercados e inteligencia aplicada al consumidor en México, a partir de las ediciones históricas del Estudio Anual de la Industria de Investigación de Mercados y Opinión Pública en México.
    \item El producto de datos se ha empleado para consolidar reportes ejecutivos de la Edición XXII (2019-2019), dirigidos tanto para el público interesado\footnote{Véase \url{ttps://www.amai.org/descargas/Estudio_Anual_AMAI_XXII_Edicion_2019_Comunicado_Medios_Junio_2020.pdf}} y otro confidencial, restringido hacia los participantes del estudio, pues contiene información de coyuntura para la industria.
    \item A partir de dicha información se desarrolló código en SQL y R, a través del cual se determinaron diferentes elementos clave a partir de la Edición XXII del estudio en cuestión, entre los cuales destacan 1) el dimensionamiento del volumen del mercado en 2019, según la facturación alcanzada en dicho periodo, 2) la cantidad de proyectos cualitativos y cuantitativos, según su tipo, 3) la cantidad de unidades realizadas para estudios cualitativos, así como en el caso de estudios cuantitativos, la cantidad de entrevistas realizado de acuerdo al tipo correspondiente, 4) la composición de los recursos humanos la industria, en todos sus segmentos, considerando la desagregación por género, 5) identificación de principales retos, cambios y temas prioritarios a futuro para la industria en el mediano plazo, 6) el ranking de empresas de acuerdo a su nivel de facturación en 2019. 
    
    \item Cabe destacar que en este punto anterior, se involucraron estimaciones de inferencia, a manera de, cálculos agregados, así como análisis de texto.
    
    \item Finalmente, se encontraron diversos cambios en la situación de este mercado, entre los que destacan: 1) un ligero crecimiento en el valor del mercado, aunque para las empresas que participaron en la Edicióon  XXII  (2018-2019)  y  la  Edicóon  XXII(2019-2020) del Estudio Anual de AMAI (0.81\%), la mayoría de ellas decrecieron, 2) en México sigue predominando la investigación de  mercado  basada  en  corte  cuantitativo (que representa 67\% de facturación anual), 3) hablando del capital humano de la industria, se observó un disminución del 61\% de los recursos humanos, sin embargo por primera ves la proporción de hombres y mujeres es 1 a 1. 
    
\end{itemize}


\section*{Recomendaciones}

\begin{itemize}

    \item Para desarrollar el estudio con conocimiento preciso del significado de las variables involucradas, es necesario consolidar documentación de las bases de datos que reúnen las respuestas de los participantes. Para avanzar en dicha dirección, se puede retomar el material presentado en la sección \ref{data:sources}.

    \item Es deseable añadir identificadores de las empresas a las bases de datos que recogen la respuesta del estudio, para encontrarse en condiciones de unir con facilidad datos de diferentes ediciones del estudio.
    
    \item Las preguntas del Estudio Anual de AMAI pueden ampliarse para considerar la circunstancias actuales pues determinan nuevos temas de conyuntura como el impacto de Covid 19, la digitalización de servicios y el uso de nuevas tecnologías.
    
\end{itemize}
